Predicting which cryptocurrencies will outperform in the following week
is fundamentally a cross-sectional ranking task, yet prior work has
applied time-series loss functions ill-suited to ranking objectives.
We demonstrate that this mismatch is the primary reason deep learning
models often underperform ordinary least squares in standard configurations.
%
Our key methodological contribution is the use of rank-percentile
normalized targets with ListNet ranking loss: instead of softmax-weighting
raw returns, we assign each asset its cross-sectional rank percentile as
the supervision signal, eliminating scale sensitivity across heterogeneous
volatility regimes.
%
We evaluate three deep learning architectures—LSTM with cross-asset
attention, a corrected Temporal Fusion Transformer (TFT), and a novel
CrossSectionalGatedNet (CS-Gated)—against six traditional baselines on
a panel of 324 weeks $\times$ 86 cryptocurrencies.
%
The best model, CS-Gated trained on on-chain features, achieves an
annualized Sharpe ratio of 2.75 (equal-weight long-short portfolio),
while LSTM with Trump social media and ETF flow features reaches 2.13.
The results show that deep learning becomes competitive once the ranking
objective is specified correctly, although different architectures remain
best suited to different signal types.
An ablation study confirms that rank normalization together with gradient
accumulation, rather than architectural novelty alone, drives the improvement.
